\documentclass[11pt,a4paper]{article}
\usepackage[margin=1in]{geometry}
\usepackage{amsmath,amssymb,amsthm}
\usepackage{algorithm}
\usepackage{algpseudocode}
\usepackage{booktabs}
\usepackage{hyperref}

\theoremstyle{definition}
\newtheorem{definition}{Definition}[section]
\newtheorem{theorem}{Theorem}[section]
\newtheorem{lemma}[theorem]{Lemma}
\newtheorem{remark}{Remark}[section]

\title{\textbf{Theoretical Analysis, Gradient Propagation Dynamics, and Residual Scaling in Transformer Normalization Placements}}
\author{Team Scaibu \\ \small Email: \texttt{jaydeep.9664920749@gmail.com} \\ \small Architecture and Core Engine Team}
\date{October 2026}

\begin{document}
\maketitle

\begin{abstract}
This document provides the rigorous mathematical specification, gradient flow derivation, and stability analysis for transformer normalization placements: Post-LN, Pre-LN, Sandwich-LN, and DeepNorm. We formulate the residual recurrence equations for each topology, analyze the gradient scale degradation in deep networks, prove why Pre-LN preserves identity gradient highways without warm-up heuristics, derive the DeepNorm parameter bounds enabling 1000-layer stability, trace a worked numerical example, and document numerical precision safeguards.
\end{abstract}

\section{Notation and Mathematical Preliminaries}

Let $x_l \in \mathbb{R}^{B \times D}$ denote the residual stream activation tensor at layer $l \in \{1, \dots, L\}$. Let $F(z) = z W$ represent a linear or non-linear sublayer transformation (e.g., self-attention or MLP block) with weight matrix $W \in \mathbb{R}^{D \times D}$.

\begin{itemize}
    \item $B \in \mathbb{N}_{\ge 1}$: Batch size or sequence length.
    \item $D \in \mathbb{N}_{\ge 1}$: Hidden feature dimension.
    \item $L \in \mathbb{N}_{\ge 1}$: Total depth of transformer layers.
    \item $\text{LN}(z; \gamma, \beta, \epsilon)$: Layer Normalization operator with gain $\gamma \in \mathbb{R}^D$ and bias $\beta \in \mathbb{R}^D$.
    \item $\alpha \in \mathbb{R}_{> 0}$: Residual scaling factor in DeepNorm ($\alpha = (2L)^{1/4}$ for encoders).
    \item $\beta_s \in \mathbb{R}_{> 0}$: Sublayer branch scaling factor in DeepNorm ($\beta_s = (8L)^{-1/4}$).
    \item $\epsilon > 0$: Denominator regularizer ($\epsilon = 10^{-5}$).
\end{itemize}

\begin{table}[h]
\centering
\caption{Summary of Normalization Placement Topologies}
\begin{tabular}{lll}
\toprule
\textbf{Topology} & \textbf{Forward Recurrence Equation} & \textbf{Gradient Characteristics} \\
\midrule
\textbf{Post-LN} & $x_{l+1} = \text{LN}(x_l + F(x_l))$ & Gradients diminish as $\mathcal{O}(1/\sqrt{L})$; requires warmup \\
\textbf{Pre-LN} & $x_{l+1} = x_l + F(\text{LN}(x_l))$ & Unimpeded residual path; gradients $\Theta(1)$ at all depths \\
\textbf{Sandwich-LN} & $x_{l+1} = x_l + \text{LN}_2(F(\text{LN}_1(x_l)))$ & Bounded sublayer output; controls explosion in vision-LLMs \\
\textbf{DeepNorm} & $x_{l+1} = \text{LN}(\alpha x_l + \beta_s F(x_l))$ & Stabilizes deep Post-LN up to 1000+ layers \\
\bottomrule
\end{tabular}
\end{table}

\section{Problem Definition and Core Concepts}

\subsection{Intuitive Meaning}
The original Transformer used Post-LN, placing normalization after the residual addition. However, repeated normalization in Post-LN scales down the gradient signal propagating to early layers, causing vanishing gradients and training divergence without careful learning-rate warmup. Pre-LN places normalization inside the sublayer branch, leaving the residual backbone untouched ($x_{l+1} = x_l + \dots$) and ensuring stable optimization across deep architectures.

\subsection{Formal Mathematical Recurrences}
\begin{align}
\text{Post-LN: } \quad x_{l+1} &= \text{LN}(x_l + F(x_l)) \\
\text{Pre-LN: } \quad x_{l+1} &= x_l + F(\text{LN}(x_l)) \\
\text{Sandwich-LN: } \quad x_{l+1} &= x_l + \text{LN}_2(F(\text{LN}_1(x_l))) \\
\text{DeepNorm: } \quad x_{l+1} &= \text{LN}(\alpha x_l + \beta_s F(x_l))
\end{align}

\section{Algorithm Specification}

\begin{algorithm}[H]
\caption{Transformer Block Forward with Normalization Placement}
\label{alg:norm_placement}
\begin{algorithmic}[1]
\Require Input tensor $X \in \mathbb{R}^{B \times D}$, sublayer weights $W \in \mathbb{R}^{D \times D}$, gain $\gamma \in \mathbb{R}^D$, bias $\beta \in \mathbb{R}^D$, placement mode $P$, coefficients $\alpha, \beta_s > 0$, regularizer $\epsilon > 0$.
\Ensure Output tensor $X_{\text{next}} \in \mathbb{R}^{B \times D}$, sublayer tensor $F_{\text{out}} \in \mathbb{R}^{B \times D}$.
\If{$P = \text{pre\_norm}$}
    \State $Z \gets \text{LN}(X; \gamma, \beta, \epsilon)$
    \State $F_{\text{out}} \gets Z W$
    \State $X_{\text{next}} \gets X + F_{\text{out}}$
\ElsIf{$P = \text{post\_norm}$}
    \State $F_{\text{out}} \gets X W$
    \State $X_{\text{next}} \gets \text{LN}(X + F_{\text{out}}; \gamma, \beta, \epsilon)$
\ElsIf{$P = \text{sandwich\_norm}$}
    \State $Z \gets \text{LN}(X; \gamma, \beta, \epsilon)$
    \State $F_{\text{raw}} \gets Z W$
    \State $F_{\text{out}} \gets \text{LN}(F_{\text{raw}}; \gamma, \beta, \epsilon)$
    \State $X_{\text{next}} \gets X + F_{\text{out}}$
\ElsIf{$P = \text{deep\_norm}$}
    \State $F_{\text{raw}} \gets X W$
    \State $F_{\text{out}} \gets \beta_s \cdot F_{\text{raw}}$
    \State $X_{\text{next}} \gets \text{LN}(\alpha X + F_{\text{out}}; \gamma, \beta, \epsilon)$
\EndIf
\State \Return $\{X_{\text{next}}, F_{\text{out}}\}$
\end{algorithmic}
\end{algorithm}

\section{Theoretical Guarantees and Gradient Dynamics}

\begin{theorem}[Identity Gradient Highway in Pre-LN]
In an $L$-layer Pre-LN network, the gradient of the loss $\mathcal{L}$ with respect to the initial representation $x_0$ satisfies:
\begin{equation}
\frac{\partial \mathcal{L}}{\partial x_0} = \frac{\partial \mathcal{L}}{\partial x_L} \left( I + \sum_{l=0}^{L-1} \frac{\partial F(\text{LN}(x_l))}{\partial x_l} \right)
\end{equation}
guaranteeing a direct, non-vanishing path from the loss directly to the input embeddings regardless of depth $L$.
\end{theorem}

\begin{proof}
By telescoping the Pre-LN recurrence $x_{l+1} = x_l + F(\text{LN}(x_l))$, the output representation satisfies:
\begin{equation}
x_L = x_0 + \sum_{l=0}^{L-1} F(\text{LN}(x_l))
\end{equation}
Differentiating with respect to $x_0$ yields $\frac{\partial x_L}{\partial x_0} = I + \sum_{l=0}^{L-1} \frac{\partial F(\text{LN}(x_l))}{\partial x_0}$. By the chain rule $\frac{\partial \mathcal{L}}{\partial x_0} = \frac{\partial \mathcal{L}}{\partial x_L} \frac{\partial x_L}{\partial x_0}$, the identity term $I$ ensures that gradients do not decay exponentially with depth.
\end{proof}

\subsection{Limitations}
\begin{itemize}
    \item \textbf{Pre-LN Output Head Norm}: Because residual variances accumulate along the depth in Pre-LN ($\text{Var}(x_L) = \mathcal{O}(L)$), a final LayerNorm before the linear projection head is strictly required to prevent token representation explosion.
\end{itemize}

\section{Complexity Analysis}

\subsection{Time Complexity}
\begin{itemize}
    \item \textbf{LayerNorm Reductions}: Takes $\mathcal{O}(B \cdot D)$ FLOPs per normalization pass.
    \item \textbf{Linear Sublayer Transformation}: Matrix multiplication $X W$ takes $2 B D^2$ FLOPs.
    \item \textbf{Total Time Complexity}: $\Theta(B \cdot D^2)$ FLOPs dominated by the matrix multiplication.
\end{itemize}

\subsection{Space Complexity}
\begin{itemize}
    \item \textbf{Output and Intermediate Tensors}: Stores $B \times D$ activations, requiring $\mathcal{O}(B \cdot D)$ space.
\end{itemize}

\section{Exactness and Error Analysis}

The placement logic executes exact floating-point arithmetic. Numerical precision matches IEEE 754 standards, with normalization subroutines employing double-precision accumulators to prevent rounding drift across multiple stacked layers.

\section{Worked Numerical Example}

Let $B = 1, D = 2$.
\begin{equation}
X = [[2.0, 4.0]], \quad W = \begin{bmatrix} 1.0 & 0.0 \\ 0.0 & 1.0 \end{bmatrix}, \quad \gamma = [1.0, 1.0]^T, \quad \beta = [0.0, 0.0]^T, \quad \epsilon = 10^{-5}
\end{equation}

\begin{enumerate}
    \item \textbf{Pre-LN Mode}:
    \begin{align}
    \mu &= 3.0, \quad \sigma^2 = \frac{(2-3)^2 + (4-3)^2}{2} = 1.0, \quad \text{std} = \sqrt{1.00001} \approx 1.000005 \\
    \text{LN}(X) &= \left[ \frac{2-3}{1.000005}, \frac{4-3}{1.000005} \right] \approx [-0.999995, 0.999995] \\
    F_{\text{out}} &= \text{LN}(X) W = [-0.999995, 0.999995] \\
    X_{\text{next}} &= X + F_{\text{out}} = [2.0 - 0.999995, 4.0 + 0.999995] = [1.000005, 4.999995]
    \end{align}
    \item \textbf{Post-LN Mode}:
    \begin{align}
    F_{\text{out}} &= X W = [2.0, 4.0] \\
    X + F_{\text{out}} &= [4.0, 8.0] \\
    \mu_{\text{sum}} &= 6.0, \quad \sigma_{\text{sum}}^2 = \frac{(4-6)^2 + (8-6)^2}{2} = 4.0, \quad \text{std} = \sqrt{4.00001} \approx 2.0000025 \\
    X_{\text{next}} &= \left[ \frac{4-6}{2.0000025}, \frac{8-6}{2.0000025} \right] \approx [-0.9999988, 0.9999988]
    \end{align}
\end{enumerate}

\section{Numerical Considerations and Edge Cases}

\begin{itemize}
    \item \textbf{Layer Freezing and Fine-Tuning}: Changing normalization placement between pre-training and fine-tuning causes instant model divergence due to incompatible internal coordinate scaling.
\end{itemize}

\begin{thebibliography}{9}
\bibitem{xiong2020layer} R.~Xiong et al., ``On Layer Normalization in the Transformer Architecture,'' in \emph{International Conference on Machine Learning (ICML)}, pp.~10524--10533, 2020.
\bibitem{wang2022deepnet} H.~Wang et al., ``DeepNet: Scaling Transformers to 1,000 Layers,'' \emph{arXiv preprint arXiv:2203.00555}, 2022.
\bibitem{vaswani2017attention} A.~Vaswani et al., ``Attention Is All You Need,'' in \emph{Advances in Neural Information Processing Systems (NeurIPS)}, pp.~5998--6008, 2017.
\end{thebibliography}

\end{document}
